% =============================================================================
% bab1-pendahuluan.tex — BAB I PENDAHULUAN
% Skripsi / Tesis / Disertasi (Lampiran I)
% =============================================================================

\chapter{PENDAHULUAN}
\label{bab:pendahuluan}

\section{Latar Belakang Penelitian}
\label{sec:latar-belakang}

Tuliskan latar belakang penelitian di sini. Uraikan fenomena, kondisi empiris, gap
penelitian, atau permasalahan yang melatarbelakangi penelitian ini. Gunakan data
dan fakta yang relevan serta didukung oleh sumber pustaka yang sahih.

Setiap paragraf diawali dengan indentasi sebesar 5 karakter (±1.27 cm) secara otomatis.
Teks menggunakan spasi 1.5 dan font Times New Roman 12pt sesuai ketentuan FEB
Universitas Telkom \cite{contohSumber2025}.

\section{Identifikasi Masalah}
\label{sec:identifikasi-masalah}

Berdasarkan latar belakang di atas, permasalahan yang dapat diidentifikasi dalam
penelitian ini adalah sebagai berikut:

\begin{enumerate}
  \item Masalah pertama yang diidentifikasi.
  \item Masalah kedua yang diidentifikasi.
  \item Masalah ketiga yang diidentifikasi.
\end{enumerate}

\section{Batasan Masalah}
\label{sec:batasan-masalah}

Penelitian ini dibatasi pada hal-hal berikut:

\begin{enumerate}
  \item Batasan pertama penelitian ini.
  \item Batasan kedua penelitian ini.
\end{enumerate}

\section{Rumusan Masalah}
\label{sec:rumusan-masalah}

Berdasarkan latar belakang dan batasan masalah di atas, rumusan masalah dalam
penelitian ini adalah:

\begin{enumerate}
  \item Bagaimana \ldots?
  \item Apakah terdapat \ldots?
  \item Seberapa besar pengaruh \ldots?
\end{enumerate}

\section{Tujuan Penelitian}
\label{sec:tujuan}

Berdasarkan rumusan masalah yang telah diuraikan, tujuan penelitian ini adalah:

\begin{enumerate}
  \item Untuk mengetahui \ldots
  \item Untuk menganalisis \ldots
  \item Untuk mengukur \ldots
\end{enumerate}

% Untuk Tesis (S2) dan Disertasi (S3), tambahkan sub-bab Manfaat Penelitian jika diperlukan
% \section{Manfaat Penelitian}
% \label{sec:manfaat}
% ...
